\section{Superseded values and withdrawn results}
\label{app:history}

This appendix is the note's audit trail. It holds the values that have been struck, the
arguments that were stated and then withdrawn, and the superseded uncertainty constructions with
their figures. The body quotes only current values and, where a passage was moved here, states
the current position in one or two sentences with a pointer to this appendix; a first reading can
skip it. Struck values are shown struck through and are not quotable. Except for headings,
joining sentences and cross-reference targets, the text of parts~B and~C and of the context
passages in part~A is reproduced verbatim from the note as it stood before this restructure,
which is commit \texttt{152ef7a6} of the MINERvA-OmniFold repository (\repolink). Part~A
also covers the struck values that commit \texttt{95d03169} (2026-08-23) removed from the body
instead of moving; their rows and sentences are taken from that commit's parent,
\texttt{95d03169\^{}}. The
repository's index of retracted and superseded values is routed in App.~\ref{app:provenance}
(\texttt{retracted-values}).

\subsection{Struck values}
\label{app:history-errata}

Each row is one struck value as the note printed it, with the reason and the record the note
gave for it. A value struck in several places is listed once, with every place it was printed.

{\footnotesize
\setlength{\tabcolsep}{3pt}
\begin{longtable}{@{}>{\raggedright\arraybackslash}p{0.23\textwidth}
                  >{\raggedright\arraybackslash}p{0.20\textwidth}
                  >{\raggedright\arraybackslash}p{0.11\textwidth}
                  >{\raggedright\arraybackslash}p{0.21\textwidth}
                  >{\raggedright\arraybackslash}p{0.20\textwidth}@{}}
\toprule
Struck & Quantity (where it was printed) & Current value & Reason & Record \\
\midrule
\endhead
$\dead{1.008\%}$ & 2D \texttt{Flux} (PPFX) per-bin median relative $\sigma$, band and grouped category (Table~\ref{tab:band-medians}; \S\ref{sec:universes}) & $4.993\%$ & Pre-flux-fix rollup: each PPFX universe divided by $\Phi_{\rm CV}$ instead of by its own flux integral $\Phi_u$ & \texttt{...\_matcorr\_fluxfix/uq\_universe\_summary.txt} \\
$\dead{4.783\%}$ & 2D systematic total per-bin median (Table~\ref{tab:band-medians}; provenance note, App.~\ref{app:history-combine}) & $6.830\%$ & Pre-flux-fix rollup, as above & \texttt{...\_matcorr\_fluxfix/uq\_universe\_summary.txt} \\
$\dead{2.463\times10^{-39}}$ & 2D systematic $\sqrt{\Tr\mathbf{C}}$ (Table~\ref{tab:band-medians}; provenance note, App.~\ref{app:history-combine}) & $3.214\times10^{-39}$ & Pre-flux-fix rollup, as above & \texttt{...\_matcorr\_fluxfix/uq\_universe\_summary.txt} \\
$\dead{2.470\times10^{-39}}$ & 2D combined (block-sum) $\sqrt{\Tr\mathbf{C}}$ (provenance note, App.~\ref{app:history-combine}) & $3.220\times10^{-39}$ & Pre-flux-fix rollup, as above & \texttt{2d-unfolding/2D\_OMNIFOLD\_STUDY\_STATUS.md}, \S``Stage-2 UQ envelope'' \\
$\dead{4.822\%}$ & 2D combined per-bin median relative $\sigma$ (provenance note, App.~\ref{app:history-combine}; worked example of \S\ref{sec:chisq}) & $6.865\%$ & Pre-flux-fix rollup, as above & \texttt{2d-unfolding/2D\_OMNIFOLD\_STUDY\_STATUS.md}, \S``Stage-2 UQ envelope'' \\
$\dead{\approx\!70\%}$ & OmniFold 2D envelope as a fraction of the paper's (provenance note, App.~\ref{app:history-combine}) & Not quoted; the envelope matches the paper ($6.865\%$ against $6.86\%$) & An inference drawn from the pre-flux-fix covariance, retracted with it & None given beyond the provenance note \\
$\dead{0.069}$ & 2D pull mean, combined covariance (\S\ref{sec:pull}; provenance note, App.~\ref{app:history-combine}) & $0.051$ & Computed on the pre-flux-fix combined covariance & \texttt{VALIDATION\_LEDGER.md} \\
$\dead{0.466}$ & 2D pull RMS, combined covariance (\S\ref{sec:pull}; provenance note, App.~\ref{app:history-combine}) & $0.409$ & Computed on the pre-flux-fix combined covariance & \texttt{VALIDATION\_LEDGER.md} \\
$\dead{0.563\%}$ & 2D grouped \emph{Statistical} category median (\S\ref{sec:universes}) & $0.549\%$ & Re-quoted from the pure-Poisson $N=300$ bootstrap, which superseded a seed-varying set that leaked ML stochasticity into the statistical block & \texttt{2d-unfolding/2D\_OMNIFOLD\_STUDY\_STATUS.md}, 2026-05-29 \\
$\dead{1.828\times10^{-40}}$ & 2D statistical $\sqrt{\Tr\mathbf{C}}$ (provenance note, App.~\ref{app:history-combine}) & $1.817\times10^{-40}$ & As above: superseded seed-varying bootstrap set & \texttt{2d-unfolding/2D\_OMNIFOLD\_STUDY\_STATUS.md}, 2026-05-29 \\
$\dead{0.564\%}$ & 2D statistical per-bin median relative spread (\S\ref{sec:bootstrap}; provenance note, App.~\ref{app:history-combine}) & $0.549\%$ & As above: superseded seed-varying bootstrap set & \texttt{2d-unfolding/2D\_OMNIFOLD\_STUDY\_STATUS.md}, 2026-05-29; current spreads \texttt{VL162} \\
$\dead{5.0467}$ & Full-event $C_{\rm stat}$ family: nominal-to-member ratio of unfolded yield in $6<p_{\parallel}<\SI{20}{GeV/c}$ (App.~\ref{app:cstatlimit}) & $3.5969$ (family mean; s.d.\ $1.6091$) & A single member's value, a $70$th-percentile draw quoted as if typical & \texttt{state/probe-oi126-band-\allowbreak Rpush-sigma-20260815} (family spread) \\
\midrule
\multicolumn{5}{@{}l}{\emph{Removed from the body on 2026-08-23 by \texttt{95d03169}; text from its parent, \texttt{95d03169\^{}}}} \\
\midrule
$\dead{\SI{2.796e-38}{cm^2/nucleon}}$ & Recoil-only PET absolute total on data, in the \texttt{niter}${}={}$2 cross-check against GBDT (removed 2026-08-23, 95d03169) & Not quoted & \texttt{niter}${}={}$2 legacy, withdrawn as a current comparison: the pinned policy is \texttt{niter}${}={}$3 (CLM-010, CLM-012), the 2026-08-01 full-event schema makes a pre-08-01 PET total a different estimator, and the missing background subtraction (J21) biases PET high & \texttt{INDEX-retracted-and-\allowbreak superseded-values.md} \\
$\dead{\SI{3.066e-38}{}}$ & GBDT absolute total on data, the other operand of that cross-check (removed 2026-08-23, 95d03169) & Not quoted & Struck with the cross-check it is an operand of, on the grounds above & \texttt{INDEX-retracted-and-\allowbreak superseded-values.md} \\
$\dead{\petRatio}$ & PET/GBDT absolute total ratio, in that cross-check and in the Fig.~\ref{fig:petabs} caption (removed 2026-08-23, 95d03169) & Not quoted & Arithmetically correct, the exact quotient of the two struck operands; withdrawn because its operands describe a training configuration that no longer exists & \texttt{INDEX-retracted-and-\allowbreak superseded-values.md} \\
\dead{\SIrange{6.5}{9.9}{\percent}} & Per-axis median PET/GBDT differences in that cross-check (removed 2026-08-23, 95d03169) & Not quoted & As above & \texttt{INDEX-retracted-and-\allowbreak superseded-values.md} \\
$\dead{\sim\SI{\petGbdtGap}{\percent}}$ & PET--GBDT data-side gap (removed 2026-08-23, 95d03169) & Not quoted & Struck beside the scope note, which covers every quantity derived from the cross-check's struck figures & \texttt{INDEX-retracted-and-\allowbreak superseded-values.md} \\
$\dead{-1.6\%}$ & Shift of the PET total under a higher-epoch retrain (removed 2026-08-23, 95d03169) & Not quoted & Struck on the same grounds, being the same legacy run & \texttt{INDEX-retracted-and-\allowbreak superseded-values.md} \\
$\dead{0.897}$ & Ratio of the PET total under that higher-epoch retrain (removed 2026-08-23, 95d03169) & Not quoted & Struck on the same grounds, being the same legacy run & \texttt{INDEX-retracted-and-\allowbreak superseded-values.md} \\
$\sim$\,\dead{\SI{\petClosure}{\percent}} & Ordinary-closure level of the recoil-only PET, in the Fig.~\ref{fig:petabs} caption (removed 2026-08-23, 95d03169) & Not quoted & Legacy two-iteration self-consistency diagnostic, not a current validation & None given in the deleted text \\
\dead{\SI{\gbdtFiveBlockMedian}{\percent}} & 5D candidate background-aware block sum: median per-bin uncertainty (removed 2026-08-23, 95d03169) & Not quoted & Pre-J28, and under the 2026-07-12 construction-cause quarantine (two independent grounds); no replacement magnitude is quoted or authorized, because the corrected J28 products are non-background-aware & \texttt{PROCEDURE-gbdtFive-\allowbreak macro-update.md} \\
$\dead{\SI{\gbdtFiveAdoptTrace}{cm^2/nucleon}}$ & 5D candidate mean-centered covariance including cross-source nonlinear response, $\sqrt{\mathrm{Tr}\,C}$ (removed 2026-08-23, 95d03169) & Not quoted & As above & \texttt{PROCEDURE-gbdtFive-\allowbreak macro-update.md} \\
$\dead{\SI{\gbdtFiveMeanShift}{cm^2/nucleon}}$ & 5D candidate joint mean-shift norm (removed 2026-08-23, 95d03169) & Not quoted & As above & \texttt{PROCEDURE-gbdtFive-\allowbreak macro-update.md} \\
$\dead{\SI{\gbdtFiveCVTrace}{cm^2/nucleon}}$ & 5D candidate CV-centered construction, the conservative variant, $\sqrt{\mathrm{Tr}\,C}$ (removed 2026-08-23, 95d03169) & Not quoted & As above & \texttt{PROCEDURE-gbdtFive-\allowbreak macro-update.md} \\
\bottomrule
\end{longtable}
}

\subsubsection{The struck entries in context}
\label{app:history-errata-context}

\paragraph{Struck values removed from the body on 2026-08-23.}
Commit \texttt{95d03169} deleted the sentences below instead of moving them; they are reproduced from its parent, \texttt{95d03169\^{}}. The unstruck prose it deleted with them, including the two scope notes that the reasons in the table above are taken from, is not reproduced.

From \S\ref{sec:pet-abs}, the third gate of the absolute path:
\textbf{Cross-check vs GBDT} (\textbf{\texttt{niter}${}={}$2 legacy --- not
        a current number}; see the scope note below): on data, the PET absolute total was
        $\dead{\SI{2.796e-38}{cm^2/nucleon}}$ versus the GBDT
        $\dead{\SI{3.066e-38}{}}$, ratio $\dead{\petRatio}$, with per-axis median
        differences of \dead{\SIrange{6.5}{9.9}{\percent}} (Fig.~\ref{fig:petabs}).

From the paragraph after its scope note:
What survives independently of the withdrawn comparison: the
$\dead{\sim\SI{\petGbdtGap}{\percent}}$ data-side gap is not
a normalization bug, but ordinary closure alone cannot assign it uniquely to
training configuration. A higher-epoch retrain moved the PET total by
$\dead{-1.6\%}$ (ratio $\dead{0.897}$) --- struck on the same grounds, being the
same legacy run --- which demonstrated training sensitivity qualitatively.

From the caption of Fig.~\ref{fig:petabs}:
The PET/GBDT total
  ratio is $\dead{\petRatio}$ --- \textbf{struck: \texttt{niter}${}={}$2 legacy, see
  the scope note in the text}; ordinary closure was assessed to
  $\sim$\,\dead{\SI{\petClosure}{\percent}} --- \textbf{struck: legacy two-iteration
  self-consistency diagnostic, not a current validation} --- and it does not test
  omitted muon dependence.

From \S\ref{sec:uq-combine}:
For the 5D candidate, the background-aware block sum gave median per-bin uncertainty
\dead{\SI{\gbdtFiveBlockMedian}{\percent}}; including cross-source nonlinear response
raised the candidate mean-centered covariance to
$\sqrt{\mathrm{Tr}\,C}=\dead{\SI{\gbdtFiveAdoptTrace}{cm^2/nucleon}}$. The joint mean
shift, norm $\dead{\SI{\gbdtFiveMeanShift}{cm^2/nucleon}}$, was reported separately
rather than folded into that covariance. A CV-centered construction gave the
larger $\dead{\SI{\gbdtFiveCVTrace}{cm^2/nucleon}}$ and was retained as a conservative
variant.

\paragraph{From \S\ref{sec:bootstrap}, the bootstrap block.}
The $N=50\to300$ comparison of the quoted replicas carried the parenthesis
(the superseded seed-varying set gave a median of $\dead{0.564\%}$), and the seed-pinning statement continued: ``It
superseded an earlier set whose replicas varied both seeds and so leaked ML
stochasticity into the statistical block (App.~\ref{app:history-combine}).'' The standing assumption on independence of error sources added
(The quoted $N=300$ bootstrap pins the GBDT seed;
        only the superseded seed-varying set leaked, \S\ref{sec:bootstrap}.)

\paragraph{From \S\ref{sec:invert}, rank deficiency at $N\le n$.}
After ``has rank at most $N-1$'' the paragraph continued: ``This was the regime of the Stage-1
bootstrap ($N=50$ replicas, $n=205$ bins): Cholesky failed without jitter, and the required
jitter, $1.6\times10^{-94}$, was a numerical artifact of the rank deficiency rather than a real
near-singularity.'' The body keeps the rank bound and the full rank of the Stage-2 ($N=300$)
bootstrap.

\paragraph{From \S\ref{sec:universes}, Table~\ref{tab:band-medians}.}
The two struck rows, as the table printed them:
\begin{center}
\begin{tabular}{@{}lc@{}}
\toprule
Band & Per-bin median rel.\ $\sigma$ \\
\midrule
Flux (PPFX, $N_{u}=100$) & $\dead{1.008\%}$ $\to$ $4.993\%$ \\
\midrule
\shortstack[l]{Systematic total\\(44 bands $+$ norm)} & \shortstack[l]{$\dead{4.783\%} \to \boxed{6.830\%}$ (median)\\$\sqrt{\Tr\mathbf{C}}=\dead{2.463\times10^{-39}} \to 3.214\times10^{-39}$} \\
\bottomrule
\end{tabular}
\end{center}
\textbf{Struck entries in Table~\ref{tab:band-medians} are the pre-flux-fix rollup}
(\texttt{uq/universe\_stage2\_MEFHC\_full\_matcorr/uq\_universe\_covariance\_full\_matcorr.root});
the live values to their right are the flux-fixed rollup
\texttt{...\_full\_matcorr\_fluxfix/uq\_universe\_summary.txt}, which is the
product this note quotes (\S\ref{sec:systematics}). Of the individual systematic
bands only \texttt{Flux} moved; the grouped \emph{Statistical} category also moved
($0.563\% \to 0.549\%$) for the separate reason given in the provenance note in
App.~\ref{app:history-combine}.
Grouped by category the medians are: Muon reconstruction $3.366\%$,
Models $2.034\%$, Flux $\dead{1.008\%} \to 4.993\%$, Hadronic response $0.758\%$,
Statistical $\dead{0.563\%} \to 0.549\%$, Normalization $0.081\%$.

\paragraph{From \S\ref{sec:chisq}, worked example.}
The median row read:
\begin{center}
\begin{tabular}{@{}lcc@{}}
\toprule
Metric (205 bins) & Paper & OmniFold (combined) \\
\midrule
Per-bin median relative $\sigma$ & $6.86\%$ & $\dead{4.822\%} \to 6.865\%$ \\
\bottomrule
\end{tabular}
\end{center}

\paragraph{From \S\ref{sec:205bins}, the strict-interior mask.}
Before its current justification the sentence on the $185$-bin mask opened: ``The mask dates from
an early stage of the analysis, when the paper's unreported diagonal cells dominated the
$\chi^{2}$;''.
% Source: pre-freeze 2D_OMNIFOLD_RUN_LOG_ARCHIVE.md:50-59 (Phase 4, unreported bins), :716-726.

\paragraph{From \S\ref{sec:pull}.}
With the combined
covariance (paper TotalCov $+$ flux-fixed matcorr universe $+$ bootstrap-300
$+$ ML) the pull mean is $\dead{0.069} \to \boxed{0.051}$ (centered on zero; no
global offset) and the pull RMS $\dead{0.466} \to \boxed{0.409}$.
The struck values were computed on the pre-flux-fix combined covariance and are
\textbf{retracted}; the live pair is the flux-fixed one recorded in
\texttt{VALIDATION\_LEDGER.md} (App.~\ref{app:history-combine} provenance note).

\paragraph{From App.~\ref{app:cstatlimit}, the measured spread and the recomputation route.}
An earlier internal statement of the band ratio used a single member's value,
$\dead{5.0467}$ --- \textbf{struck: a $70$th-percentile draw quoted as if typical}; the
family-typical value is $3.5969$ with the spread above.
That receipt carries two retracted share-of-total fields, marked in place;
neither relative-standard-deviation field is among them, and no share-of-total figure from it is used
here.

\subsection{Withdrawn arguments}
\label{app:history-withdrawn}

Each passage below stated an inference and withdrew it in place. The body keeps only the
surviving statement.

\paragraph{From \S\ref{sec:invert}, pseudo-inverse.}
After giving NumPy's documented \texttt{pinv} cutoff and the looser $n\,\sigma_{\max}\varepsilon_{\text{machine}}$ convention, the sentence ended
``\ldots; an earlier
version of this appendix quoted that formula for \texttt{pinv} and was wrong.''
The comparison with the looser convention read: ``\ldots\ convention of \texttt{np.linalg.lstsq}
and MATLAB, which at $n=205$ would be about $45\times$ looser
($205\times2.22\times10^{-16}=4.6\times10^{-14}$)''. The body keeps the operative cutoff
($\texttt{rcond}=10^{-15}$) and the statement that it is not that convention.

\subsubsection{From \S\ref{sec:combine}: the provenance note on the 2D budget}
\label{app:history-combine}
\begin{quote}
\textbf{Provenance note (2026-07-31; table corrected 2026-08-11).} Table~\ref{tab:combine}
carries the \emph{flux-fixed} rollup. An earlier draft printed the
\emph{pre-flux-fix} product --- systematic $\dead{2.463\times10^{-39}}$ /
$\dead{4.783\%}$, combined $\dead{2.470\times10^{-39}}$ / $\dead{4.822\%}$ ---
and inferred from it that the OmniFold envelope is $\dead{\approx\!70\%}$ of the
paper's. Those four figures and that inference are \textbf{retracted}: they are
artefacts of a superseded covariance, not of the measurement. The current budget
agrees with \S\ref{sec:systematics}, where the envelope essentially
\emph{matches} the paper ($6.865\%$ against $6.86\%$).

The two rollups differ in \emph{exactly one band}. Comparing
\texttt{2d-unfolding/uq/universe\_stage2\_MEFHC\_full\_matcorr/} against
\texttt{...\_matcorr\_fluxfix/}, every systematic band agrees to the last quoted
digit except \texttt{full\_Flux}, whose per-bin median rises
$1.008\% \rightarrow 4.993\%$ once each PPFX universe is divided by its own flux
integral $\Phi_u$ instead of by $\Phi_{\rm CV}$. That single correction carries the
universe total $4.783\% \rightarrow 6.830\%$ and the combined figure to
$6.845\%$ before the ML block, which is the whole of the discrepancy. The pull
widths of \S\ref{sec:pull} were affected identically:
$\dead{0.069}/\dead{0.466}$ were computed on the same superseded covariance, and
the flux-fixed values are $0.051/0.409$ (\texttt{VALIDATION\_LEDGER.md}).

The statistical cell moved for an unrelated and earlier reason, recorded because
it is a correction in its own right and \emph{not} because it changed the
headline: it is quoted from the pure-Poisson $N=300$ bootstrap
($\dead{1.828\times10^{-40}}$ / $\dead{0.564\%}$ $\to$
$1.817\times10^{-40}$ / $0.549\%$), which superseded a seed-varying set that leaked
ML stochasticity into the statistical block
(\texttt{2d-unfolding/2D\_OMNIFOLD\_STUDY\_STATUS.md}, 2026-05-29).
\textbf{It is not what moved the combined row.} The block sum rounds to
$3.220\times10^{-39}$ with either statistical cell --- the change sits about two
orders of magnitude below the printed precision --- and essentially the whole
movement of the combined entry comes from the systematic cell, i.e.\ from the flux
fix. This is stated explicitly so that a later check does not reverse a correct
edit on finding that the headline is insensitive to it.

The same defect, uncorrected, is still present in the N-D and 5-D kernels; see
\texttt{AUDIT-FINDINGS-20260731.md} J18 and J28.
\end{quote}

\paragraph{From \S\ref{sec:chisq}, limit (ii) of the worked example.}
The heading of the limit, and the end of its paragraph, read:

\emph{(ii) ``Correct uncertainties give exactly one'' is false and has been
removed from this appendix.}

An earlier version
of this paragraph said that a perfectly correct model with correct
uncertainties gives exactly $1$, and then used $\approx1.5$ as a tolerance
against it. Both halves of that inference are withdrawn.

\paragraph{From \S\ref{sec:closure-hv}, hidden-variable closure.}
After ``$A=0.30$ gives $3.43\%$, which \emph{exceeds} $3.0\%$'':
\S\ref{sec:validation}
records that arm as exceeding the threshold and this appendix agrees with it;
the earlier reading here --- that $A=0.30$ was under threshold and the leakage
was ``small enough to be safely ignored'' --- was an arithmetic error and is
withdrawn.

\paragraph{From App.~\ref{app:cstatlimit}, construction.}
An
earlier version reported from them that MC accounted for ${\sim}88\%$ of the variance and the
background MC alone for $87\%$; that step does not follow and is withdrawn.

\paragraph{From App.~\ref{app:cstatlimit}, the two readings of the band.}
After ``The family disagrees with the central value by a factor of order three, coherently, in the best-accepted region of the plane.'':

Two readings are available a priori.
\textbf{(a) The spread is real:} the estimator is genuinely this sensitive to the measured
sample there, and the quoted uncertainties in $6<p_{\parallel}<\SI{20}{GeV/c}$ really are of
order $67\%$. \textbf{(b) The proxy is wrong:} $\mathrm{Poisson}(1)$ resampling gives $36.8\%$
of measured rows zero weight, and a learned estimator facing that much removed support may
respond in a way that is a property of the perturbation rather than of the data. Both fit the
measurements above: an unstable estimator and an invalid perturbation both act during training,
concentrate where information is scarcest, and leave the extraction faithful.
The proposed binary was withdrawn on 2026-08-19 as a type mismatch: a covariance measures
dispersion \emph{about a centre} and does not establish finite-family \emph{containment} of a
separately trained nominal, so non-containment of the nominal by $50$ draws is not what this
covariance reports and \textbf{no choice between (a) and (b) follows from it}. The readings are
recorded because the measurements were gathered while they were live and a reader will meet them in
that form elsewhere; neither is asserted.

The resampling argument that follows opened with the sentence
``What follows concerns the resampling scheme alone and is offered as an argument rather than a
ratified result; it bears on reading (b) without adjudicating the retracted fork.''
Within it:
An earlier version concluded that the scheme ``is the sampling
distribution rather than a stand-in that might be the wrong one''; that promotion is withdrawn.

In the paragraph on the discriminating $\mathrm{Exponential}(1)$ contrast, which the body keeps, the
sentence before the surviving objection read, with the opening of that objection:
No ``publish it as it stands''
branch exists for either outcome to select; the earlier framing here --- that an unimportant zero
atom would mean publishing the covariance as it stands --- described a posture the ruling had already
closed, and is withdrawn. What stands on its own is the other half: \ldots

In ``What the band figure is'':
An earlier version called it an ``upper bound on how poorly the data constrain the
cross section,'' reasoning that a better-regularised estimator could improve on it and none could do
worse for want of information. That is withdrawn: no information bound follows from an observed
refit dispersion without conditions on the estimator's bias and on the perturbation's validity as a
sampling proxy, and neither is established here; a more strongly regularised estimator can show
\emph{less} dispersion while constraining the cross section no better, exactly the direction that
breaks the claimed inequality.

\paragraph{From App.~\ref{app:cstatlimit}, the resampling argument.}
The subsection on the band, then titled ``Why the two readings of the band are not a live fork'',
opened: ``The family disagrees with the central value by a factor of order three, coherently, in
the best-accepted region of the plane. Two readings were proposed as a binary --- \textbf{(a)} the
spread is real, \textbf{(b)} the $\mathrm{Poisson}(1)$ proxy is wrong --- and no choice between
them follows from this covariance, which measures dispersion about a centre and does not establish
finite-family containment of a separately trained nominal.'' The resampling argument that followed
read, in full:

What follows concerns the resampling scheme alone and is offered as an argument rather than a
ratified result; it bears on reading (b). The resampled
objects are individual reconstructed events. A $\mathrm{Poisson}(1)$ multiplicity $k$ on a row
states that, were the data retaken, that event occurred $k$ times; zero states that it did not
occur, which is exactly what counting noise does to an event observed once, so the zero atom is not
an artifact of the method. Independent $\mathrm{Poisson}(1)$ multiplicities also reproduce exact
Poisson fluctuations in every aggregate of those rows (\S\ref{app:poisson}), which variance-matched
continuous alternatives do not: they match the first two moments and get the discreteness wrong.
\textbf{That argument reaches counts, and it stops there.} Exact aggregation concerns unweighted sums
of resampled rows; this covariance disperses the output of a weighted, learned, iteratively refit
estimator evaluated on every replica, and nothing above shows that a perturbation with the right
counting behaviour induces the right \emph{sampling} distribution for that functional --- still less
that the resulting interval covers.

The body keeps the surviving conclusion: the zero atom is what counting noise does to a row
observed once, so its presence is not by itself evidence that the proxy is invalid.

\paragraph{From App.~\ref{app:cstatlimit}, limitations of the covariance object.}
\textbf{The $257$ are the quotable all-member domain, and they are not an
``invertible sub-block'' --- this appendix called them one, which is impossible at this rank, and the
name is corrected here.}

\paragraph{From \S\ref{sec:coverage}, the 2D coverage test.}
The
attempted 100/100 split-sample reinterpretation is also withdrawn because it
used the all-toy mean of the fluctuated truth histograms.

\paragraph{From App.~\ref{app:negweight}, the extended-fiducial footing and the measured size of the choice.}
\textbf{The $\sigma_{\mathrm{ext}}=\SI{\fpsSigma}{cm^2/nucleon}$ reported in
\S\ref{sec:fps} was not produced with negative-weight injection}, and an earlier
version of this appendix implied that it was.

What is withdrawn is the \emph{footing}: negative-weight injection
is the intended baseline for the extended-fiducial analysis, not the one behind
the number quoted here.

It is not a theorem that the two
footings give the same covariance, because the underlying estimands are not
pointwise identical --- so ``the two covariances agree by construction,'' as
this paragraph previously read, is withdrawn.

\paragraph{From App.~\ref{sec:negweight-estimator}, the continuous estimand.}
After stating that the purity weight and negative-weight injection agree exactly only in the
binned limit and target the bin-integrated estimand $r^{\star}(b)=(D_{b}-B_{b})/S_{b}$, the
sentence continued ``\ldots; the claim that they give the same continuous estimand
$r^{\star}(x)=(D(x)-B(x))/S(x)$ was stated here without those conditions and is corrected.''

\paragraph{From \S\ref{sec:eavailw-construction}, conditioning.}
Two properties of this
matrix are routinely misstated, including in an earlier version of this
subsection, so both are given with their operands.

\paragraph{From \S\ref{sec:pet-abs}, the PET--GBDT level of agreement.}
The only such comparison that exists was taken with a two-iteration training,
against the campaign's pinned three-iteration policy (\texttt{CLM-010},
\texttt{CLM-012}, \textsc{frozen}); before the 2026-08-01 full-event schema change,
which makes a pre-08-01 PET total a different estimator rather than a
differently-trained one; and with unit measured weights and \emph{no} background
subtraction (audit finding J21), which at this analysis's
${\sim}\SI{3}{\percent}$ background scale biases PET high, so the underlying gap is
\emph{larger} than the withdrawn figure. It has not been re-extracted under the
pinned configuration, and the same three grounds withdraw the legacy closure level
and the higher-epoch retrain response.

The figure that followed that discussion in \S\ref{sec:pet-abs} plots the withdrawn comparison, and
is kept here with it.

\begin{figure}[H]
  \centering
  \includegraphics[width=0.95\textwidth]{pet_vs_gbdt_absolute}
  \caption{Recoil-only PET representation cross-check: absolutely-normalized
  4D cross section from the point-cloud unfold (full-statistics reweight of the generator sample) versus the
  production scalar GBDT result, projected on each axis. The PET/GBDT total
  offset shown is a two-iteration legacy result, withdrawn as a current comparison
  (see text) and not quoted here as a level of agreement. Ordinary closure was
  assessed only as a legacy two-iteration self-consistency diagnostic; no current
  closure level is quoted, and it does not test omitted muon dependence.}
  \label{fig:petabs}
\end{figure}

\paragraph{From \S\ref{sec:summary}, outlook.}
After listing the remaining paper-independent cross-checks, the outlook continued ``The
truncated-spectral `ours-only' goodness-of-fit convention (App.~\ref{sec:rank}) is no longer
among them: the collaboration confirmed it uses the same truncated-spectral pseudo-inverse.''

\subsection{Superseded uncertainty constructions and their figures}
\label{app:history-superseded}

None of the constructions below is a publication uncertainty, and no number in this part is
quotable.

\subsubsection{The historical 3D generator scan and uncertainty budget}
From \S\ref{sec:3d-models}, after the central-value generator ratios:

The historical truncated-spectral scan
ranked Tune~v1 closest and GiBUU farthest, and still does when redrawn on the
flux-repaired predictions, but its covariance is superseded by
the corrected dependency contract: the final 3D covariance must be projected
from the adopted, selection-complete 5D trunk. Consequently no 3D generator
$\chi^{2}$, $p$-value or significance is quoted here. The old scan is retained
only as a quarantined numerical diagnostic in Fig.~\ref{fig:3dfullcov}.

\begin{figure}[H]
  \centering
  \includegraphics[width=0.9\textwidth]{compare_3d_fullcov}
  \caption{Quarantined historical full-3D truncated-spectral scan. It shows the
  per-generator $\chi^{2}/\mathrm{ndf}$ against the historical covariance versus
  eigenvalue tolerance (redrawn 2026-09-27 on the flux-repaired GENIE and NuWro
  predictions, \S\ref{sec:3d-models}) and is
  retained only to document numerical behavior; no value is interpreted until
  the selection-complete 5D covariance is adopted \emph{and} projected to 3D.
  The first has happened, and on 2026-09-22 so did the second: the projection
  exists (\texttt{20c16e16}\ldots, 1431 reported cells). \textbf{The gate still
  holds, on different ground} --- that product is \emph{constructed, not
  adopted}, and it exists only in the receiving-cells destination-mask variant,
  not the \texttt{declared-dst-cv} variant a figure would require.}
  \label{fig:3dfullcov}
\end{figure}

\paragraph{The historical 3D uncertainty budget, and why it is not quoted.}
\label{sec:3d-syst}
The historical band on Fig.~\ref{fig:3dmodels-superseded} was built as a full 3D
analog of the 2D uncertainty budget of \S\ref{sec:systematics}. The
systematic universes are dumped on the three-axis
$(\pT,\pPar,\Eavail)$ grid (187 universes), the bootstrap
statistical and ML training-seed-variation covariances propagated, and the three summed into a
single historical combined
covariance $C_{\mathrm{syst}}+C_{\mathrm{stat}}+C_{\mathrm{ML}}$ over all 1431
reported bins. It has $\sqrt{\operatorname{tr}}=\SI{5.72e-39}{}$, median
per-bin relative uncertainty \SI{10.4}{\percent}, numerical rank $247$ (of
1431; eigenvalues above $10^{-12}\lambda_{\max}$), and the same flux-dominated band ordering as 2D
(Fig.~\ref{fig:3dbands}). These are audit descriptors, not publication
uncertainties. Under the corrected contract, the quotable 3D covariance is the
exact projection of the adopted 5D covariance. That covariance now exists
(\S\ref{sec:eavailw-construction}) and the replacement it waited on has landed, but
the 3D projection was built on 2026-09-22 (\texttt{20c16e16}\ldots).
Covariance-dependent 3D comparisons remain gated, on the two grounds that
survive: that projection is \emph{constructed, not adopted}, and it exists only
in the receiving-cells destination-mask variant, not the
\texttt{declared-dst-cv} variant a comparison against the frozen 3D central
value would require.

\begin{figure}[H]
  \centering
  \includegraphics[width=0.49\textwidth]{uq_universe_3d_band_eavail}
  \includegraphics[width=0.49\textwidth]{uq_universe_3d_band_pt}
  \caption{Grouped fractional systematic bands of the historical 3D covariance
  projected onto $\Eavail$ (left) and $\pT$ (right): flux-dominated, with the
  same band ordering as the 2D budget. The magnitudes are quarantined pending
  the final selection-complete 5D-to-3D projection.}
  \label{fig:3dbands}
\end{figure}

\paragraph{From \S\ref{sec:uq-ml}, the split-inclusive ML band on the 3D budget.}
``\ldots\ a \emph{train/test-split} variant, which also re-draws the internal training
split, gives a band $1.24\times$ the pure-seed band, moving the combined 3D
budget by only $+0.04\%$ in $\sqrt{\mathrm{Tr}\,C}$ \ldots''

\paragraph{From \S\ref{sec:3d-2p2h}, the Valencia 2p2h fraction.}
An earlier \SI{27}{\percent} came mostly from comparing two
different generated samples for CV and CV$+$MEC; the MEC component alone gave
\SI{64}{\percent} before the repair. Historical 3D $\chi^{2}$ values
are not quoted pending the corrected covariance projection.

\subsubsection{The historical recoil-PET uncertainty assembly}
\label{app:history-petbudget}
From \S\ref{sec:pet-syst}:

The historical recoil-PET campaign produced vertical, retraining, ML,
statistical and detector blocks on a common nominal.  That assembly is now
quarantined rather than quoted as an uncertainty.  For a nuisance endpoint
$u$, the physical shift and the induced retraining response form one joint
shift,
\begin{equation}
  \delta_u=x_u^{\rm varied+retrained}-x_{\rm CV}.
\end{equation}
Writing this as a frozen-map piece plus a retraining increment is an algebraic
decomposition, but adding the two separate covariances omits
$\operatorname{Cov}(\delta_{\rm frozen},\delta_{\rm retrain})$ and its
transpose.  Positive semidefiniteness of that sum does not restore the missing
cross term.

The historical detector block is additionally CV-support-limited: it propagates
shifted weights and observables through the nominal cloud membership instead of
regenerating the selected cloud for every lateral endpoint.  A replacement
construction would need end-to-end joint nuisance/retraining shifts and the
selection-complete FPS lateral samples, and would belong to a full-event PET
estimator with its own statistical and ML ensembles.

\subsubsection{The earlier direct $(\Eavail,W)$ block sum}
\label{sec:eavailw-cov}
From \S\ref{sec:eavailw}:

\noindent\textbf{This subsection is historical. It describes a superseded
construction and the requirements imposed on its replacement; the covariance
this note actually quotes is specified in \S\ref{sec:eavailw-construction}.}

An earlier attempt built a dedicated 42-bin $(\Eavail,W)$ covariance
\emph{directly}, as $C = C_{\mathrm{syst}}+C_{\mathrm{stat}}+C_{\mathrm{lateral}}$,
by a \emph{frozen-reweighter} block sum on the merged five-axis
\texttt{\_universes\_full} sample: the central-value OmniFold push-weights
were held fixed and only the per-event truth reweight ratios changed per
systematic universe (the same frozen-reweighter \emph{technique} as the PET
budget, App.~\ref{app:history-petbudget} --- not the PET covariance matrix, which is never
applied to this GBDT central value; no per-universe re-inference). That object
carried 13 knob bands $\pm1\sigma$, the 100-universe flux multisim, the
projected full five-axis bootstrap covariance $C_{EW}=M C_{5D}M^{\top}$, and a
$W$-resolved detector (lateral) block. It is not adopted and is not quoted
anywhere in this note.

The paragraph that followed, stating the three requirements declared for a replacement and that the adopted object meets them, remains in the body in the ``Construction'' paragraph of \S\ref{sec:eavailw-construction}.

\subsubsection{The superseded 5D candidate unified-throw constructions}
\label{app:history-5dcandidates}
From \S\ref{sec:uq-higherdim}:

For the 5D candidate this construction was carried out both mean-centered and
CV-centered, with the joint mean shift reported separately rather than folded into
the covariance. \textbf{No magnitude from either construction is quoted here.} Both are
pre-J28 and background-aware, both sit under the 2026-07-12 construction-cause
quarantine, and the corrected J28 products are footed on the
\emph{non}-background-aware configuration --- so no substitution is available, and
supplying one would make the sentence false rather than current. Neither is adopted for publication, and that
is no longer a statement about waiting: the selection-complete lateral
replacement has landed and the adopted product is a different object
(\S\ref{sec:eavailw-adopted}), so these two are superseded rather than
pending. No
positive-semidefiniteness measurement exists for either: \texttt{VL16} and \texttt{VL17}
record that property for the \emph{successor} background-aware products only, and it does
not transfer by inheritance.

\subsubsection{Disclosures shortened in the body}
Until 2026-10-03 the body figures below drew uncertainties from superseded covariances, and their captions said so in one clause. The fuller wording they carried is kept here. The body now shows them as central values only; the versions with those uncertainties are reproduced below with the captions they carried in the body (Figs.~\ref{fig:3dmodels-superseded}--\ref{fig:ascencio-superseded}).

\paragraph{From \S\ref{sec:systematics}.}
The 3D and 4D budgets
(\S\ref{sec:3d-syst}, \S\ref{sec:4d}) repeat the identical procedure on the
higher-dimensional grids.

\paragraph{From \S\ref{sec:3d-models}.}
Figure~\ref{fig:3dmodels} overlays the four on the unfolded result, with the
combined-covariance systematic band of \S\ref{sec:3d-syst} on the data.

\paragraph{From Fig.~\ref{fig:3dmodels}, caption.}
\textbf{The grey band labelled ``total syst$\oplus$stat'' inside the graphic
  is the SUPERSEDED historical 3D covariance} of \S\ref{sec:3d-syst}
  (\texttt{hCov\_combined3d\_total}), drawn for orientation only. It is not the
  uncertainty of this measurement, it is not quotable, and no comparison in this
  section uses it; the quotable 3D covariance must be projected from the adopted
  5D trunk; that projection now exists (2026-09-22) but is \emph{not adopted},
  so nothing here is quotable from it.

\begin{figure}[H]
  \centering
  \includegraphics[width=0.99\textwidth]{generators_vs_unfolded_band}
  \caption{1D projections of the unfolded 3D central cross section (black)
  versus four generator predictions
  produced at truth level: GENIE 2.12 CV, MINERvA Tune v1, NuWro 21.09, and
  GiBUU 2019 (GENIE and NuWro: the flux-repaired predictions,
  \S\ref{sec:3d-models}). On $\pT$ and $\pPar$ the models track the data shape; on $\Eavail$
  they split, and all under-predict the data across the quasielastic--$\Delta$
  dip ($0.1\le\Eavail<\SI{0.4}{GeV}$). The wide
  $\Eavail$ catch bin is omitted from the axis.
  \textbf{The grey band labelled ``total syst$\oplus$stat'' inside the graphic
  is the superseded historical 3D covariance}, drawn for orientation only and not
  quotable (App.~\ref{app:history}). Only the central values are interpreted.}
  \label{fig:3dmodels-superseded}
\end{figure}

\paragraph{From Fig.~\ref{fig:modedecomp}, caption.}
\textbf{The error bars drawn on the data points
    (``$\pm$ total'') are from the superseded historical 3D covariance}
    (\S\ref{sec:3d-syst}); they orient the eye and are not quotable, and the
    excess is read here as a central-value shape statement.

\paragraph{From Fig.~\ref{fig:mec}, caption.}
\textbf{As in Fig.~\ref{fig:modedecomp}, the ``$\pm$ total'' bars on the
    data points come from the superseded historical 3D covariance}
    (\S\ref{sec:3d-syst}) and are not quotable; the gap-filling fraction is a
    central-value statement.

\begin{figure}[H]
  \centering
  \begin{minipage}{0.49\textwidth}
    \centering
    \includegraphics[width=\textwidth]{mode_decomp_eavail}
    \caption{Central-value $d\sigma/d\Eavail$ decomposed by interaction mode and
    overlaid with the unfolded data. The data excess is localized to the
    quasielastic--$\Delta$ dip ($\Eavail\le\SI{0.4}{GeV}$), the shape and size of
    a missing 2p2h component. The ``$\pm$ total'' error bars are from the
    superseded historical 3D covariance and are not quotable (App.~\ref{app:history}).}
    \label{fig:modedecomp-superseded}
  \end{minipage}\hfill
  \begin{minipage}{0.49\textwidth}
    \centering
    \includegraphics[width=\textwidth]{compare_mec_eavail}
    \caption{Re-generating GENIE with the empirical Valencia 2p2h enabled (red,
    stacked on the no-2p2h CV; both flux-repaired) fills part of the low-$\Eavail$ dip --- about half
    the data$-$CV gap (\SI{52}{\percent}), the remainder being the region the
    MINERvA tune's empirical enhancement is built to absorb.
    The ``$\pm$ total'' bars are from the superseded historical 3D covariance and
    are not quotable (App.~\ref{app:history}).}
    \label{fig:mec-superseded}
  \end{minipage}
\end{figure}

\paragraph{From Fig.~\ref{fig:ascencio}, caption.}
The
  uncertainty-derived annotations inside the panel --- the error bars on our
  points and the $\chi^{2}$ it prints --- are computed from the
  \emph{superseded} historical unified-throw 4D covariance, not from a
  corrected one, and are not quotable; the published points carry the
  published covariance.

\begin{figure}[H]
  \centering
  \includegraphics[width=0.85\textwidth]{ascencio_fullcov_compare}
  \caption{Bin-identical comparison to Ascencio
  \emph{et al.}~\cite{MINERvA:2022incl} on the maximal common
  $(\Eavail,q_{3})$ grid. Only the central-value comparison is used. The
  error bars on our points and the printed $\chi^{2}$ are from the superseded
  historical unified-throw 4D covariance and are not quotable
  (App.~\ref{app:history}); the published points carry the published covariance.}
  \label{fig:ascencio-superseded}
\end{figure}

\subsubsection{Pre-repair values}
\label{app:history-prerepair}
Values from before a repair or correction, which the body mentioned in passing beside the current value it still prints.

\paragraph{From \S\ref{sec:3d-2p2h}, the high-$\Eavail$ bin.}
2p2h also does not touch the high-$\Eavail$ $[1.50,3.00)$ bin, where the data
exceed the flux-repaired GENIE CV by about \SI{9}{\percent} (data/CV $1.09$;
it was $1.15$ before the repair) \ldots

\paragraph{From \S\ref{sec:eavailw-band}, the generator inputs.}
The 3D and $(\Eavail,W)$ comparisons
now take each generator from one five-dimensional histogram, so the earlier
unreconciled \SI{3}{\percent} difference between two GENIE~CV productions no
longer enters: the in-phase-space GENIE~CV total is \SI{2.76e-38}{cm^2/nucleon}
in both. The GENIE$+$MEC prediction uses the 3D normalization convention, which
also removes a \SI{2.87}{\percent} under-normalization of the earlier
$(\Eavail,W)$ MEC prediction (\texttt{KNOWN\_ISSUES.md} row 81).
% Source-only comment block moved verbatim from sec_eavailw.tex (sec:eavailw-band),
% organization review 2026-10-04 (lane 2 sec. 8.3 finding 2). Not rendered.
% --- SUPERSEDED 2026-09-27 (s5p, OI-193; KNOWN_ISSUES 83 and 81) ---------------
% Every GENIE/NuWro number in the provenance block below (1.535/1.579/1.563, the
% 2.4446e-38 vs 2.52e-38 discrepancy, the W-band 1.30-1.35) rests on the
% flux-defective predictions and is HISTORY. GiBUU moved only by its x0.9966
% flux-convention reweight (1.609 -> 1.607). The current numbers are
% VALIDATION_LEDGER VL156-VL160, from
% docs/orchestration/state/s5p/stage7/generator-context/generator-context-receipt.json
% (eavailW_band, W_2p2_3p0) and eavailW_band_s5p_repaired-log.txt beside it:
%   corner data/gen  GENIE-CV 1.142  GENIE+MEC 1.139  NuWro 1.160  GiBUU 1.607
%   integrated       GENIE-CV 1.110  GENIE+MEC 1.068  NuWro 1.153  GiBUU 1.386
%   W in [2.2,3.0)   GENIE-CV 1.123  GENIE+MEC 1.111  NuWro 1.130  GiBUU 1.343
%   corner/integrated (derived): 1.029, 1.066, 1.006, 1.159
% The block below is kept unedited as the record of the 2026-08 reduction.
% --- PROVENANCE / REDUCTION for the corner ratios (J27.3, recorded 2026-08-02) ---
% The ratios below are CORNER-INTEGRATED cross-section ratios -- not per-bin means,
% not a weighted mean, not the worst bin. Reduction, from
% 3d-unfolding/genie/overlay_eavailW_band.py:88-108:
%   corner = outer(EAVAIL_EDGES[:-1] >= 0.8, W_EDGES[:-1] >= 1.8)  -> 3x3 = 9 cells
%            [Eavail 0.8-1.5, 1.5-3.0, 3.0-100] x [W 1.8-2.2, 2.2-3.0, 3.0-100]
%   data/gen = sum_corner(hData2D * dEavail * dW)
%              / sum_corner(hXSec_eavailW * dEavail * dW)
% Source of 1.54/1.58/1.56: ND_OMNIFOLD_RUN_LOG.md:988-990 (2026-06-08 band run),
% which reads "All THREE underpredict ... 1.54 CV, 1.58 +MEC, 1.56 NuWro" and states
% "GiBUU excluded (FinalEvents.dat lacks per-event Enu)". GiBUU was structurally
% ABSENT from the run these three numbers come from. This paragraph previously read
% "All four underpredict ... by 54-58%" -- a later edit that attached three
% magnitudes to four generators without a recomputation; corrected 2026-08-02 to
% attribute them to the three that produced them. GiBUU was afterwards added to
% the band (gibuu_to_xsec_eavailW.py, Task 16; make_figures.sh:55 passes it and the
% figure legend shows it), so the 4-generator overlay DID print a GiBUU corner
% ratio -- to stdout, never captured. To close: re-run make_figures.sh:55 and read
% the "hiE-hiW corner ... data/gen=" line. Inputs live only on Perlmutter /pscratch.
% Do NOT substitute either of these -- both are different reductions:
%   - Eavail-projected GiBUU data/gen 1.59 / 1.36 / 1.91 (RUN_LOG:1126-1127)
%     integrates over ALL W, not W>=1.8.
%   - corner chi2 381.1/12 (RUN_LOG:1155-1157) uses Eavail>=0.4 -> 12 cells,
%     a different sub-block that adds the 0.4-0.8 Eavail row.
% Until recomputed, 54-58% is a THREE-generator span. When the GiBUU ratio arrives
% this is a one-line upgrade: add it, and widen the span only if it lands outside.
%
% CLOSED 2026-08-11. All FOUR ratios recomputed TOGETHER on one data file, so the set
% is internally consistent rather than a fourth number appended to three older ones.
% Log (whole stream, no tail): 3d-unfolding/genie/eavailW_band_20260811_allfour.log
% Ledger: "2026-08-11 (E_avail,W) four-generator corner ratios". Predeclared at
% docs/orchestration/PREDECLARE-20260811-gibuu-corner-ratio.md, branch G1.
%   GENIE-CV  corner 8.7918e-39  data/gen 1.535    (note had 1.54)
%   GENIE+MEC corner 8.5484e-39  data/gen 1.579    (note had 1.58)
%   NuWro     corner 8.6369e-39  data/gen 1.563    (note had 1.56)
%   GiBUU     corner 8.3893e-39  data/gen 1.609    <- NEW; lands OUTSIDE 1.54-1.58,
%   data      corner 1.3497e-38                       so the span widens 54-58 -> 54-61
% The three reproduce AT THE PRECISION THIS NOTE QUOTES. Identity is NOT established:
% the only surviving record of the 2026-06-08 values is RUN_LOG:988-990 at three
% significant figures, and the data file (excess_eavail_W.root, 2026-07-14) POSTDATES
% them by five weeks -- so a shift in the third decimal is masked by rounding. Per
% BEN-086, agreement at printed precision is not identity; do not re-quote these at
% more digits from the 2026-06-08 lineage.
% Controls that reproduce EXACTLY and are why the set is trustworthy: GiBUU integrated
% 2.2227e-38 and data 3.0699e-38 (this paragraph's own 2.22 / 3.07). The ORDERING
% GiBUU < NuWro < GENIE CV matches sec_3d.tex:215-217, but the GENIE value does not:
% this band's GENIE-CV 2.4446e-38 (work_seed* gevgen, gen_to_xsec_eavailW.py) is a
% different object from sec_3d's 2.52e-38 (Stage-A work_p* gevgen, genie_to_xsec3d.py);
% the 2.95% difference is not reconciled. (The ledger's VL35 prose repeats the old claim.)
%
% NORMALIZATION INCONSISTENCY, FOUND WHILE CLOSING THIS AND DELIBERATELY LEFT ALONE.
% This paragraph states two deficits, two sentences apart, in TWO DIFFERENT
% NORMALIZATIONS, both rendered as a bare percentage with no marker:
%   "underpredict ... by 54-61%"  is generator-relative, (data-gen)/gen = data/gen - 1
%   "sit 23-26% below the data"   is data-relative,      (data-gen)/data
% Both are internally correct and each matches its own arithmetic. But in a COMMON
% normalization the corner deficit is 34.9-37.8% data-relative (or the high-W deficit
% is 29.8-34.8% generator-relative), so a reader comparing "54-61%" against "23-26%"
% infers a contrast of ~2.4x where the consistent answer is ~1.2x. Choosing the
% convention is an authorial decision about a physics claim, so it was ROUTED rather
% than silently changed; the ratios above are unaffected either way.
%
% RESOLVED 2026-08-12 by Joseph (-> Session A -> Session B, BEN-082(v)): "use one
% normalization throughout -- prefer explicit data/generator ratios for both
% comparisons rather than mixing 'percent above generator' with 'percent below data.'"
% Both sentences now state data/generator and neither states a bare percentage.
% The W-band ratios were COMPUTED FROM THE SOURCE CROSS-SECTIONS, not converted from
% the rounded 23-26%, because converting rounded endpoints would manufacture precision
% the source does not have (the BEN-086 trap this same comment block warns about).
% From eavailW_band_20260811_allfour.log:22-23, W in [2.2,3.0) GeV, units cm^2/nucleon:
%   data 7.481e-39
%   GENIE-CV  5.764e-39 -> 1.2979    GENIE+MEC 5.622e-39 -> 1.3307
%   NuWro     5.715e-39 -> 1.3090    GiBUU     5.550e-39 -> 1.3479
%   span 1.298-1.348 -> quoted 1.30-1.35
% CROSS-CHECK, and it is why these are trustworthy: the same operands reproduce the
% note's previous data-relative wording exactly -- 22.95%, 24.85%, 23.61%, 25.81%,
% i.e. "23-26% below the data". The prose changed convention; no number moved.
% The corrected contrast is now legible on the page: 1.54-1.61 against 1.30-1.35.
% --------------------------------------------------------------------------------

\paragraph{From \S\ref{sec:3d-models} and \S\ref{sec:eavailw-band}, the superseded ledger rows.}
The generator comparisons of \S\ref{sec:3d-models} and \S\ref{sec:eavailw-band} were
rebuilt on 2026-09-27 by re-running the unchanged figure producers on the repaired inputs
(App.~\ref{app:provenance}, \texttt{flux-repair}). Ledger rows \texttt{VL156}--\texttt{VL160}
(\texttt{VL156}--\texttt{VL158} for the $(\Eavail,W)$ band) supersede the pre-repair rows
\texttt{VL35}--\texttt{VL38}.

\paragraph{From \S\ref{sec:uq-higherdim}, the background-CV choice.}
``\ldots the effect on the \emph{adopted}
covariance is smaller still, \SI{0.09}{\percent} before the flux (J28)
correction and \SI{0.18}{\percent} after it \ldots''
The fuller wording read:
``Repeating all 187 universe unfolds and the CV with
per-universe background subtraction changes the combined \emph{block-sum}
sqrt-trace by only \SI{0.28}{\percent}; the effect on the \emph{adopted}
covariance is smaller still, \SI{0.18}{\percent} after the flux (J28)
correction, because the adoption's per-bin
$\max$ inflation transfer (Eq.~\eqref{eq:zinflation},
\S\ref{sec:eavailw-construction}) damps the block-sum change. These are ratios taken against
the superseded products [the 5D candidates of \S\ref{app:history-5dcandidates}], so they bound the size of the background-CV
choice rather than describing a current covariance; either way that concern is
numerically negligible.''

\paragraph{From \S\ref{sec:pet-projection}, cloud coverage before the miss-row fix.}
Originally the cloud was filled only in the reco-signal event loop, so the
truth-only ``miss'' events appended for the unfold (\SI{\pcMiss}{\percent} of
the truth sample) carried no cloud --- only \SI{\pcHasCloudBefore}{\percent} of
truth events had one, and cloud projections were correspondingly
acceptance-limited, landing exactly on the stored-$\Eavail$ projection
restricted to the has-cloud subset but sitting $1.5$--$2\times$ below
the full unfold per bin. This is a truth-side instance of the same
dangling-miss-row dump issue tracked for the systematic bank. We fixed it at
source by carrying the truth final-state hadrons through the
truth-denominator loop and filling \texttt{part\_gen} for the appended miss
rows (the reco cloud \texttt{part\_reco} correctly stays empty for a miss),
then re-dumping, rebuilding the inputs, and re-training.

\paragraph{From \S\ref{sec:pet-projection}, the \texttt{ISSUE-30} projection fix.}
After ``the unfolded counts are already acceptance-corrected'':
``\ldots so
that division corrected for acceptance twice. The fix lowers every projected
$d\sigma/d\Eavail$ bin by 25--51\% and leaves the spread above essentially unchanged.''
